\begin{definition}
    Вероятностная мера $\pi$ на $X$ называется \textit{стационарной} или \textit{инвариантной меры цепи}, если $\pi = \pi\P$.
\end{definition}


В случае конечного $X$ такая мера \textbf{всегда существует}. 

В общем случае инвариантная мера \textbf{не единственна}, но в следующей теореме, называемой \textit{эргодической}, выделен важный случай, когда она единственна, причем при любом начальном распределении распределения $\pi_n$ сходятся к инвариантному. В данном конечном случае все разумные виды сходимости
мер совпадают, так что можно говорить о сходимости по вариации (по
метрике $\|\mu - \nu\| =
\P|\mu(x_i) - \nu(x_i)|)$ или на каждом элементе из $X$.

\begin{theorem}[Условие эргодичнности конечной цепи Маркова]
    Пусть $X$ конечно и при некотором $n_1$ все элементы матрицы $\P^{n_1}$ положительны. Тогда инвариантная вероятностная мера $\mu$ цепи единственна, причем к ней сходятся меры $\pi_n = \pi_0\P^n$ для всякого начального распределения цепи и $\mu(x_i) > 0$ для всех $i$.
\end{theorem}

\begin{note}
    Отметим, что условие положительности некоторой степени $P$ и необходимо для справедливости заключения теоремы. В самом деле, если всякая матрица $P^n$ имеет нулевой элемент, то для некоторых $i_0, j_0$ для бесконечной последовательности $n_k$ матрица $P^{n_k}$ имеет нулевой элемент с индексами $i_0, j_0$. Тогда для меры $\pi$, равной нулю на всех $x_i$ c $i \neq i_0$ и $\pi(x_{i_0}) = 1$, получаем $\pi P^{n_k} (x_{j_0}) = 0$, поэтому предел не может быть положительным на всех $x_j$.
\end{note}